\section{Evaluation}\label{sec:eval}

We evaluate \tool{} against the tools a ZK project would use today. Four
research questions guide the evaluation:
\begin{description}
\item[RQ2] Detection: does \tool{} find known binding flaws that existing
tools miss?
\item[RQ3] New flaws: does it find binding flaws that auditors missed?
\item[RQ4] Confirmation: how many warnings reach executable evidence?
\item[RQ5] Cost: what does a scan cost in money, time and human effort?
\end{description}

\subsection{Setup}\label{sec:eval-setup}

\paragraph{Development and test split}
We split the labeled findings by project family, never by item, so that no
project appears in both sets. The split is seeded and was fixed before anyone
read the test code. Families seen by the authors before the split were forced
into development.
% src: dataset/labels-v1/build_v1.py (SEED 20261124; FORCED_DEV)
\tool{} was built and tuned on the development split only, and frozen before
any test run.

\paragraph{Benchmarks}
We use two benchmarks.
\begin{itemize}
\item \textbf{B1, audit test set.} These are the audited vulnerable
commits of the test families, with the flaw located in code by the
code-level labeling pass.
  \begin{itemize}
  \item \emph{Round~1}: 14 binding flaws, 6 hard negatives (proof-related
  findings that are not binding, at their own locations) and 6 fixed
  versions.
  \item \emph{Round~2}: 8 binding flaws and 8 hard negatives located by a
  later code-level pass whose results the tool-building process could not
  read. Round~2 is our clean retest.
  \end{itemize}
\item \textbf{B2, binding mutants.} 27 binding flaws injected into three
small hosts (a privacy vault, a red-packet claim, and an oracle claim, one of
them a proof of concept) that belong to no dataset family. There are at least
four per pattern (P1 4, P2 8, P3 5, P4 5, P5 5), and 4 of them are
circuit-side edits. Each mutant is designed to keep the honest path
accepting. All 27 compile, and for one host the existing test suite confirms
it: the clean host passes 32 of 32 tests, and each mutant fails only the
tests of the removed check. The unmodified host is the clean control.
\TODO{E9}{honest-path check for the other two hosts}
\end{itemize}
% src: mvp/eval/GT-locations.csv; benchmarks/binding-mutants-v2/README.md
All targets are cleaned snapshots: the version history, audit reports and
audit references are removed, so no tool can read the answer.

\paragraph{Tools}
We compare against the tools a ZK project would use today.
\begin{itemize}
\item Static Solidity analyzers: Slither~\cite{slither} and
Mythril~\cite{mythril}.
\item Circuit analyzers: circomspect~\cite{circomspect} and
Picus~\cite{picus}.
\item Property-based fuzzing: Foundry invariant testing~\cite{foundry}.
\item The same LLM used as a plain auditor, without skills or tools.
\end{itemize}
Every tool sees the same snapshots. We give the static tools build-enabled
copies, with dependencies installed at the pinned versions, because most
snapshots do not compile without them. We run the fuzzer in two modes:
\emph{out of the box}, with generic properties, and \emph{assisted}, with the
generic binding invariant of Section~\ref{sec:design-validator} wired by hand.
The assisted mode shows what a fuzzer can do once someone supplies the
specification that \tool{} infers.

\paragraph{Scoring}
The matching rules were committed before each tool ran. One amendment (v1.1)
was adopted after the first round's results were seen, and we disclose it. It
adds a case-insensitive matcher and the separate mechanism-level metric, and
it applies to all tools equally.
\begin{itemize}
\item \emph{Strict} hit: the warning names the flaw's file and function, and
an LLM judge confirms the same mechanism.
\item \emph{Loose} hit: the location alone matches.
\item \emph{Mechanism-level} hit: the right mechanism is reported anywhere in
the repository.
\end{itemize}
A target that a tool cannot build or analyze is \emph{unsupported}. It is
reported separately and never counted as a miss. Intervals are 95\% Wilson
intervals.
% src: BindCheck/dataset/labels-v1/SCORING.md; mvp/eval/SCORING-v1.1.md

\subsection{RQ2: Detection}\label{sec:eval-rq2}

\begin{table}[t]
\centering
\caption{What each tool can see. \checkmark{} = analyzes the layer;
$\circ$ = partial; -- = not at all.}
\label{tab:capability}
\footnotesize
\setlength{\tabcolsep}{3.5pt}
\begin{tabular}{@{}lcccc@{}}
\toprule
\textbf{Tool} & \textbf{Circuit} & \textbf{Verifier} & \textbf{Consumer} & \textbf{Proof$\to$effect} \\
\midrule
Slither, Mythril & -- & \checkmark & \checkmark & -- \\
circomspect, Picus & \checkmark & -- & -- & -- \\
Foundry (out of box) & -- & $\circ$ & \checkmark & -- \\
Foundry (assisted) & -- & mock & \checkmark & manual \\
Plain LLM & $\circ$ & $\circ$ & \checkmark & $\circ$ \\
\tool{} & \checkmark & \checkmark & \checkmark & \checkmark \\
\bottomrule
\end{tabular}
\end{table}

\begin{table*}[t]
\centering
\caption{Detection on the audit test set B1 (strict hits / positives;
95\% Wilson interval). ``Sup.'' = targets the tool could analyze, out of
26 in round~1. Specificity counts fixed versions and hard negatives without a
warning at the flaw's location. Red cells are pending runs.}
\label{tab:b1}
\footnotesize
\begin{tabular}{@{}lccccc@{}}
\toprule
\textbf{Tool} & \textbf{Sup.\ (R1)} & \textbf{Round 1 strict} & \textbf{Round 2 strict} & \textbf{Fixed-version spec.} & \textbf{Hard-neg.\ spec.} \\
\midrule
Slither & 15/26 & 0/14 (0.00--0.22) & \TBD & 6/6 & 3/6 \\
Mythril & 13/26 & 0/14 (0.00--0.22) & \TBD & 6/6 & 6/6 \\
circomspect & 1/26 & 0/14 (0.00--0.22) & \TBD & 6/6 & 5/6 \\
Picus & 0/26 & unsupported & unsupported & -- & -- \\
Foundry, out of box$^\dagger$ & 13 & \multicolumn{2}{c}{0/13 (R1+R2)} & -- & -- \\
Foundry, assisted$^\dagger$ & 13 & \multicolumn{2}{c}{5/13 (0.18--0.64)} & 4/5 & 3/3 \\
Plain LLM (same model) & 26/26 & 3/14 (0.08--0.48) & \TBD & 6/6 & 6/6 \\
\tool{} & 26/26 & \DEMO{8/14 (0.33--0.79)} & \DEMO{6/8 (0.41--0.93)} & \DEMO{6/6} & \DEMO{13/14} \\
\bottomrule
\multicolumn{6}{@{}l}{\scriptsize $^\dagger$Fuzzing covers EVM-side flaws of both rounds; circuit- and guest-side flaws are unsupported. Harnesses were wired by hand in 5--45 min per target.}
\end{tabular}
% src: mvp/eval/COMPARISON.md (slither-build, mythril-build, circomspect-build, baseline-default); mvp/eval/RESULTS-static.md (supported counts); mvp/eval/fuzz/RESULTS-fuzz.md
\end{table*}
\TODO{E1}{BindCheck v2 frozen run on B1 round 1 + round 2 and B2 (in progress); replace red cells}
\TODO{E8}{static tools and the plain LLM on the round-2 targets, for complete columns}

\paragraph{Existing tools miss binding flaws}
Table~\ref{tab:b1} shows the result on real audited flaws. Neither Solidity
analyzer reports any of the 14 round-1 flaws at the right location with the
right mechanism. Slither and Mythril report something at the flaw's location
for 3 and 2 of the 14 flaws, but the judge finds that none of these reports
describes the binding mechanism. The circuit analyzers fail for a different reason: most test
flaws do not live in a circuit at all. circomspect could analyze one target,
and Picus could not compile any. Out of the box, the fuzzer found none of the
13 flaws it could run on, despite up to millions of calls per target. Its generic properties
(no revert, no panic) are never violated, because a binding flaw does not
crash anything: the contract works exactly as written.
% src: mvp/eval/RESULTS-static.md; fuzz/RESULTS-fuzz.md

\paragraph{The missing piece is the specification}
Given our binding invariant, the same fuzzer finds 5 of the 13 flaws and raises no alarm on 3 of 3 hard negatives. One fixed version
is flagged. Whether this is a false positive or a remaining gap in that fix is
still open. The
difference between 0 and 5 is entirely the specification. That
specification, which effect values must follow from the statement, is
exactly what \tool{}'s ledgers make explicit. Writing it by hand took 5--45
minutes per target and could not express circuit-side or zkVM-guest flaws.
% src: fuzz/RESULTS-fuzz.md (assisted HITs L035, L015, L010, L042, L193; L024 fixed flagged)

\paragraph{\tool{} finds what the others miss}
\DEMO{\tool{} finds 8 of the 14 round-1 flaws and 6 of the 8 round-2 flaws at
the exact location and mechanism, against 3 of 14 for the same model without
the skill procedure. It does so without a single warning on any fixed
version.} Round~2 matters most here, because the tool-building process
never had access to these flaws or their locations.
\DEMO{On the 6 positives that no other tool found and that are in scope for
the fuzzer, \tool{} finds 4.}

\begin{table}[t]
\centering
\caption{Detection of injected binding mutants (B2) by pattern
(strict hits / mutants). ``uns.'' = the tool cannot analyze the layer or
the run failed. Red cells are pending.}
\label{tab:b2}
\footnotesize
\begin{tabular}{@{}lcccccc@{}}
\toprule
\textbf{Tool} & \textbf{P1} & \textbf{P2} & \textbf{P3} & \textbf{P4} & \textbf{P5} & \textbf{All} \\
 & (4) & (8) & (5) & (5) & (5) & (27) \\
\midrule
Slither & 0 & 0 & 0 & 0 & 0 & 0 (12 uns.) \\
Mythril & 0 & 0 & 0 & 0 & 0 & 0 (4 uns.) \\
circomspect & -- & 1$^*$ & -- & -- & -- & 1 \\
Picus & -- & 0 & -- & -- & -- & 0 \\
\tool{} & \DEMO{4} & \DEMO{6} & \DEMO{4} & \DEMO{5} & \DEMO{3} & \DEMO{22} \\
\bottomrule
\multicolumn{7}{@{}l}{\scriptsize $^*$a style warning (\texttt{<--}) at the edited line, not a binding diagnosis.}
\end{tabular}
% src: benchmarks/binding-mutants-v2/RESULTS-static.md
\end{table}

\paragraph{Mutants}
B2 tests every pattern on hosts outside the dataset (Table~\ref{tab:b2}).
Slither and Mythril again report no binding flaw at the mechanism level:
Mythril returns zero issues on all 23 consumer mutants, and Slither fails to
analyze one of the hosts. The only circuit-tool hit is a style warning that
happens to sit on the edited line.
% src: binding-mutants-v2/RESULTS-static.md
\DEMO{\tool{} detects 22 of the 27 mutants and reports no binding warning on
any of the three clean hosts. Its misses are two P5 mutants that rely on
truncating a 256-bit value, two circuit-only P2 mutants, where the brain did
not connect the edited constraint to the consumer, and one P3 mutant whose
replay key is derived in a library.}

\finding{Existing analyzers find none of the binding flaws, and a fuzzer finds
them only when a human supplies the binding specification. \tool{} infers
that specification and \DEMO{finds 6 of 8 held-out audited flaws and 22 of 27
injected ones}.}

\subsection{Why \tool{} misses}\label{sec:eval-misses}

We inspected every positive that \tool{} missed, on development and test
data. Three causes recur. They are confirmed on development data and
\DEMO{account for 4 of the 6 test misses}.
\begin{itemize}
\item \emph{Parser differentials.} The proof covers one parse of the bytes,
and the application acts on another. In the vlayer email case, the DKIM-signed
region and the body that the application reads are parsed by different code.
No ledger row is wrong, and the flaw lies inside a parser.
\item \emph{SDK-level validation.} The missing check sits in a client SDK
method that the consumer trusts, outside the analyzed contract (an SP1 log
method without anchor resolution).
\item \emph{Deep dependencies.} A replay key or encoding is derived in a
library that is not vendored, so the brain cannot read it.
\end{itemize}
The first two causes need analyses beyond the proof-to-effect link. The third
is an engineering limit of our syntactic tools (Section~\ref{sec:disc-limits}).
% src: mvp/tool/FREEZE-v2.md known weaknesses 3 (L003 parser differential, L008 SDK method)
\TODO{E1}{classify the test misses into these causes}

\paragraph{Case study}
\DEMO{In a round-2 proof marketplace, \tool{}'s effect ledger lists the
fulfillment entry with the value \emph{request identifier} and the verdict
\emph{needs statement}. The statement ledger shows that the proof's journal
does not commit to that identifier. The final warning names the function, the
unbound identifier, and the attack: anyone who observes a valid fulfillment
can attach it to a different request and collect that request's payment.
Neither analyzer reports anything at that location, and the assisted fuzzer
confirms the same flaw only after a 30-minute manual harness.}
\TODO{E1}{replace with a real round-2 detection once results are in}

\subsection{RQ3: New binding flaws}\label{sec:eval-rq3}

\DEMO{We ran \tool{} on the latest audited commits of 25 ZK applications
from our collection that were not used in development, and on the current
versions of the same projects. It reported 41 warnings, of which we confirmed
11 as binding flaws by manual review. Three are previously unreported, in two
projects, and all three reach E4. In one, a recovery module accepts a proof
over an account's guardians but takes the new owner set and threshold from
calldata, so any party that observes a valid recovery proof can install its
own owners. Table~\ref{tab:new} lists them. We reported all three to the
maintainers.}
\TODO{E3}{hunt: run frozen BindCheck on the hunt pool, triage, E4-confirm, disclose; the author decides the live-version scope}

\begin{table}[t]
\centering
\caption{Previously unreported binding flaws found by \tool{}.
Projects are anonymized until disclosure completes.}
\label{tab:new}
\footnotesize
\begin{tabular}{@{}llllc@{}}
\toprule
\textbf{ID} & \textbf{App type} & \textbf{Pattern} & \textbf{Effect} & \textbf{Evidence} \\
\midrule
\DEMO{N1} & \DEMO{account recovery} & \DEMO{P1} & \DEMO{owner takeover} & \DEMO{E4} \\
\DEMO{N2} & \DEMO{\TBD} & \DEMO{\TBD} & \DEMO{\TBD} & \DEMO{\TBD} \\
\DEMO{N3} & \DEMO{\TBD} & \DEMO{\TBD} & \DEMO{\TBD} & \DEMO{\TBD} \\
\bottomrule
\end{tabular}
\end{table}

\subsection{RQ4: Confirmation}\label{sec:eval-rq4}

\DEMO{Of the 39 test-set warnings that matched a known flaw or were confirmed
by manual review, 31 (79\%) reached E2 with an automatically generated
harness, and 18 (46\%) reached E4. Every E4 attempt on a false positive failed
to commit an unauthorized effect, so no false positive was upgraded. The
warnings that stop at E1 are mostly circuit-side or zkVM-guest flaws, for
which we cannot yet generate an attacker transaction automatically.}

\begin{table}[t]
\centering
\caption{Evidence reached by confirmed warnings (RQ4). Red: pending.}
\label{tab:evidence}
\footnotesize
\begin{tabular}{@{}lrrrr@{}}
\toprule
\textbf{Benchmark} & \textbf{Confirmed} & \textbf{E1 only} & \textbf{E2} & \textbf{E4} \\
\midrule
B1 round 1 & \DEMO{8} & \DEMO{3} & \DEMO{5} & \DEMO{3} \\
B1 round 2 & \DEMO{6} & \DEMO{2} & \DEMO{4} & \DEMO{3} \\
B2 mutants & \DEMO{22} & \DEMO{3} & \DEMO{19} & \DEMO{12} \\
New flaws (RQ3) & \DEMO{3} & \DEMO{0} & \DEMO{3} & \DEMO{3} \\
\midrule
Total & \DEMO{39} & \DEMO{8} & \DEMO{31} & \DEMO{18}\textsuperscript{a} \\
\bottomrule
\multicolumn{5}{@{}l}{\scriptsize \textsuperscript{a}E4 needs the project's prover; E2 counts include E4.}
\end{tabular}
\end{table}
\TODO{E5}{run the validator (A2, A3) on all confirmed test warnings and report the E1/E2/E4 shares}

\subsection{RQ5: Cost}\label{sec:eval-rq5}

On the development split, a scan costs \$1.05 on average and takes a median
of 267\,s, with a maximum of about 11 minutes. Large monorepos cost \$2--3.
% src: mvp/tool/FREEZE-v2.md
\DEMO{On the test set, the median is 290\,s and \$1.10 per target.} The
static analyzers run in seconds to minutes, but only after a human has made
each project build. The assisted fuzzer needs 5--45 minutes of expert wiring
per target before its first run.
% src: fuzz/RESULTS-fuzz.md (wiring minutes)
\TODO{E1}{test-set cost and time from the v2 run}

\subsection{Variance and contamination}\label{sec:eval-validity}

\paragraph{Run-to-run variance}
LLM output varies between runs.
\DEMO{Across three independent runs on round~2, \tool{} finds 5, 6 and 6 of
8 flaws strictly; on B2, 21, 22 and 22 of 27.} We report the median.
\TODO{E2}{three full runs on B1 and B2; report min/median/max}

\paragraph{Training-data contamination}
The audit reports are public and may be in the model's training data. Three
results argue against memorization.
\begin{itemize}
\item The same model without the procedure finds only 3 of 14 round-1 flaws,
so the knowledge alone does not produce the detections.
\item B2 mutants are new code that no model has seen.
\item \DEMO{Detection on findings published after the model's training cutoff
is not lower than on older findings (5/7 vs.\ 9/15).}
\end{itemize}
\TODO{E6}{contamination split: detection by report date vs.\ model cutoff}
